\documentclass[11pt,letterpaper]{article}
\usepackage{technical-user-manual}
\setDocTitle{DevSecOps and Supply-Chain Assurance Reference}
\setDocSubtitle{Engineering technical reference}
\setDocVersion{1.0}
\setDocStatus{Technical education / proposed reference}
\setDocOwner{CI/CD Engineering}
\setDocDate{30 October 2026}
\setDocSummary{Supplementary engineering practice}
\begin{document}
\makedoctitle
\makedocsummary
\tableofcontents
\clearpage
\section{Purpose, scope and authority}
This is an original modern engineering reference. The 2007 Duvall, Matyas and Glover book informs foundational CI intent but does not prescribe GitHub-specific implementations or 2026 supply-chain standards.

\section{Control pipeline}
Source review and branch rules; secrets prevention and scanning; dependency analysis and license policy; SAST; tests and coverage; artifact scanning; SBOM; provenance attestations; independent verification; protected promotion and runtime monitoring.

Verification: collect the configuration or runtime evidence supporting this control, identify an accountable owner, and record the exception path where the condition cannot presently be met.

\section{Threat boundaries}
A PR contributor may control source and workflow content but must not gain a deployment token. Third-party actions execute code on the runner. Build artifacts can be substituted between stages. Self-hosted runners can persist cross-job state if not reset.

Verification: collect the configuration or runtime evidence supporting this control, identify an accountable owner, and record the exception path where the condition cannot presently be met.

\section{Triaging vulnerabilities}
A scanner finding is evidence to evaluate. Record affected dependency or dataflow, exploitability, exposure, severity, ownership, due date and exception rationale. Avoid using raw severity as the sole release criterion.

Verification: collect the configuration or runtime evidence supporting this control, identify an accountable owner, and record the exception path where the condition cannot presently be met.

\section{Attestation and SLSA}
SLSA v1.2 has distinct Source and Build tracks; individual controls do not automatically confer a level. Build provenance binds subject digests to verified build information. Verify attestations and artifact identity before deployment rather than merely generating a file.

Verification: collect the configuration or runtime evidence supporting this control, identify an accountable owner, and record the exception path where the condition cannot presently be met.

\section{Gate policy}
Critical confirmed reachable exposures block release unless a formally approved compensating control and expiration exist. False positives require reproducible reasoning and a reviewer. Secret exposure triggers rotation, not only deletion of a commit.

Verification: collect the configuration or runtime evidence supporting this control, identify an accountable owner, and record the exception path where the condition cannot presently be met.

\section{Assurance evidence}
Capture SHA, trusted workflow identity, scan config versions, scanner results, signed/attested artifact digest, reviewer, policy version and disposition. Avoid publishing sensitive SARIF, logs or private source details.

Verification: collect the configuration or runtime evidence supporting this control, identify an accountable owner, and record the exception path where the condition cannot presently be met.

\section{Adoption checklist}
\begin{itemize}
\item Identify applicable repositories, application owners and trust boundaries.
\item Define minimum pass criteria and measurable acceptance thresholds.
\item Review workflow permissions, external actions, outputs and retention.
\item Record exceptions with justification, owner and expiry.
\item Revalidate after action, runner or policy changes.
\end{itemize}
\end{document}
